\section*{ЗАКЛЮЧЕНИЕ}
\phantomsection
\addcontentsline{toc}{section}{Заключение}

В ходе выполнения курсовой работы проведён численный анализ влияния дискретного микрорельефа с \varDimpleTypeGenitive{} профилем углублений на трибологические характеристики радиального подшипника скольжения. В рамках работы решены следующие задачи:

\begin{enumerate}
    \item Выполнен обзор литературы по гидродинамической теории смазки и исследованиям лазерного текстурирования поверхностей трения. Показано, что основным механизмом улучшения характеристик подшипника является микрогидродинамический эффект, возникающий на кромках углублений за счёт формирования локальных клиновых зазоров.

    \item Изучена методика численного решения уравнения Рейнольдса для подшипника конечной длины на основе метода конечных разностей с итерационной схемой Гаусса-Зейделя и верхней релаксацией. Вычисления выполнены на сетке $500 \times 500$ узлов при критерии сходимости $\delta < 10^{-5}$.

    \item Рассчитаны поля давления и интегральные характеристики (нагрузочная способность $F$, коэффициент трения $\mu$, расход смазки $Q$) для гладкого подшипника и подшипника с \varDimpleTypeGenitive{} углублениями (тип~3) в диапазоне $\varepsilon = 0{,}05 - 0{,}80$. Геометрические параметры подшипника: $R = 35$~мм, $c = 50$~мкм, $L = 56$~мм.

    \item Проведено сопоставление результатов для гладкого и текстурированного подшипников. При рабочем эксцентриситете $\varepsilon = 0{,}6$ получены следующие результаты:
    \begin{itemize}
        \item нагрузочная способность возросла на $39{,}24$\%: с $F = 5028{,}5$~Н до $F = 7001{,}5$~Н;
        \item коэффициент трения уменьшился на $30{,}03$\%: с $\mu = 0{,}0154$ до $\mu = 0{,}0108$;
        \item расход смазки увеличился на $1{,}97$\%: с $Q = 0{,}2359$~л/с до $Q = 0{,}2405$~л/с.
    \end{itemize}
\end{enumerate}

Результаты расчётов свидетельствуют о том, что нанесение цилиндрических углублений с параметрами $a = 2{,}4$~мм, $b = 2{,}2$~мм, $h_p = 10$~мкм ($h_p/c = 0{,}2$) приводит к значительному улучшению трибологических характеристик подшипника. Одновременный рост нагрузочной способности и снижение коэффициента трения при умеренном увеличении расхода смазки подтверждают целесообразность применения данного типа микрорельефа.

Цилиндрический профиль углублений характеризуется резкими кромками и ступенчатым изменением зазора, что формирует интенсивные локальные градиенты давления. Это обеспечивает более выраженный микрогидродинамический эффект по сравнению с плавными профилями, однако может приводить к повышенным локальным напряжениям в смазочном слое. Для подшипников, эксплуатируемых в стационарных режимах при средних и высоких эксцентриситетах, цилиндрический микрорельеф является эффективным решением.

По итогам проведённого исследования можно рекомендовать использование \varDimpleTypeGenitive{} микрорельефа для подшипников скольжения нефтегазового оборудования, функционирующих в режиме жидкостного трения. Наибольшую практическую значимость текстурирование приобретает для тяжелонагруженных узлов ($\varepsilon > 0{,}5$), где прирост нагрузочной способности и снижение потерь на трение достигают максимальных значений.
